\section{Implementation details}\label{app:impl}

This appendix expands \cref{sec:methods,sec:protocol}: \cref{tab:methods} expands \cref{tab:knobs}, the paragraphs below give the pipeline, the MUSIQ input, source versions and every arm's formula behind \cref{alg:sega,alg:ctm}, and \cref{app:laws} gives the per-band laws.

\begin{table}[!htbp]
\centering
\caption{\textbf{Arms, their knobs and sample sizes (expands \cref{tab:knobs}).} Temperature is the logit gain at $4096^2$; t/x: varies with the noise level / with the image. Held-out arms ran on rows 8--194, the pilot and dev arms only on the eight dev prompts. The shaded row is the reference, \segaofficial. Key result: \tempsb{} matches \segaofficial{} in base and temperature and drops only the per-band gain.}
\label{tab:methods}
\footnotesize
\setlength{\tabcolsep}{2.6pt}
\begin{tabular}{@{}l l c c c l@{}}
\toprule
Arm & Frequency base & Temperature & Per-band gain & t/x & $n$ \\
\midrule
\ntk & $\theta\,s^{2D/(D-2)}$ & none & none & no / no & \N{ho@nprompts} \\
\dype & PI/NTK ramps in $\sigma$ & $\N{c@dypegain}$ (spatial) & none & yes / no & \N{ho@nprompts} \\
\rowcolor{papertint}\segaofficial & $\theta\,s^{2D/(D-2)}/(1{+}0.1\ln s)$ & $s^{0.16}=\N{c@segatemp}$ & sensor, $\ge1$ & yes / yes & \N{ho@nprompts} \\
\segapaper & $\theta\,s^{2D/(D-2)}$ & $s^{0.16}$ & sensor, unbounded & yes / yes & 8 (pilot) \\
\tempsb & as \segaofficial & $s^{0.16}$ on all dims & none & no / no & \N{ho@nprompts} \\
\constancy & as \ntk & $\tau_N=\N{c@taulogn}$ & gated, late target & yes / yes & \hon{cons} (partial) \\
\constancytm & as \segaofficial & $s^{0.16}$ & gated, time-matched & yes / yes & \N{ho@nprompts} \\
\aliastemp & as \segaofficial & none & alias count, static & no / no & 8 (dev) \\
\segactm & as \segaofficial & $s^{0.16}$ & \sega{} $\times$ time-matched & yes / yes & \N{ho@nprompts} \\
\bottomrule
\end{tabular}
\end{table}

\begin{table}[!htbp]
\centering
\caption{\textbf{Configuration shared by every arm.} Gray bands group model, sampling and evaluation settings from the pipeline description below; \cref{tab:methods} lists what differs between arms. Key result: the arms differ only in rotary frequencies and gains.}
\label{tab:config}
\footnotesize
\setlength{\tabcolsep}{4pt}
\begin{tabular}{@{}l l@{}}
\toprule
\rowcolor{black!7}\multicolumn{2}{@{}l}{\textit{Model}} \\
Transformer & \krea{}, 19 double-stream and 38 single-stream blocks \\
Attention & 24 heads of dimension 128; rotary blocks of 16, 56 and 56 dimensions \\
Text encoders and autoencoder & T5 (512 tokens) and CLIP; \fluxdev{} autoencoder with tiled decoding \\
\rowcolor{black!7}\multicolumn{2}{@{}l}{\textit{Sampling}} \\
Resolution & $4096\times4096$, a $\N{c@tokside}^2$ token grid \\
Sampler & \N{c@steps} Euler flow-matching steps, flow shift $\mu=\N{c@muofficial}$, distilled guidance 4.5 \\
Precision and batch & bfloat16, batch size one, per-prompt CPU seed \\
\rowcolor{black!7}\multicolumn{2}{@{}l}{\textit{Evaluation}} \\
Prompts & \afk{} rows 0--7 (dev) and 8--194 (held-out, $n=\N{ho@nprompts}$) \\
Seeds & 0 for every arm; 1 on rows 8--62 for three arms \\
Statistics & paired percentile bootstrap, 2,000 resamples, generator seed 0 \\
Software & diffusers 0.39, PyTorch 2.11, pyiqa 0.1.16 \\
\bottomrule
\end{tabular}
\end{table}

\paragraph{Pipeline}
The \krea{} transformer, loaded through an environment override (\texttt{CROPE\_\allowbreak TRANSFORMER}), runs inside the \fluxdev{} pipeline (T5 and CLIP text encoders with 512 T5 tokens, autoencoder with tiled decoding at $4096^2$), diffusers 0.39~\citep{vonplaten2022diffusers}, PyTorch 2.11, bfloat16, batch size one, a per-prompt CPU seed, \N{c@steps} Euler steps and distilled guidance 4.5. The transformer has 19 double-stream and 38 single-stream blocks with 24 heads of dimension 128, and normalizes q and k before the rotation. Every arm multiplies q and k of every token in every block by one gain vector after the rotation; the vector is shared by blocks and heads and recomputed once per model call. The scheduler patch that holds $\mu=\N{c@muofficial}$ is applied before any configuration is recorded. The per-image JSON records hard-code the base model name and an attention label; the loaded transformer comes from the task scripts and the attention policy from the code revision (\cref{sec:pitfalls,app:compute}).

\paragraph{MUSIQ input pipeline}
We use pyiqa 0.1.16's KonIQ10k MUSIQ model. The project passes RGB float32 values in $[0,1]$ without resizing; MUSIQ normalizes to $[-1,1]$ and combines the native image with aspect-preserving 224- and 384-pixel views, resized by bicubic interpolation with \texttt{align\_corners=False} and \texttt{antialias=False}. Each scale is split into $32\times32$ patches at stride 32, with SAME padding where needed. The native setting \texttt{max\_seq\_len\_from\_original\_res=-1} retains all patches: a $4096^2$ frame contributes 16,384 native patches, plus 49 and 144 resized-view patches, for 16,577 before the class token. A $1024^2$ crop contributes 1,217. Neither native square grid needs padding. The spatial embedding's $10\times10$ hash grid does not downsample the image. We verified these counts by executing the installed preprocessor without loading model weights. Crop MUSIQ averages four separately scored, fixed random $1024^2$ crops shared across methods. Full-frame MUSIQ is thus a multiscale quality score that includes every native pixel, not a direct sharpness measurement.

\paragraph{Source versions}
Our \sega{} comparison uses \href{https://github.com/rajabi2001/sega/tree/ddaab08552c7b7ceba2ebd8f579a0f65dee90b5e}{released commit \texttt{ddaab08552c7}}. The allocation is in \texttt{flux\_sega/transformer\_flux.py}, lines 103--154; projected-feature sensing follows the input projection in the same file. Our printed-law implementation follows Section 4.1 and equations (3)--(7) of \citet{sega2026}, with the strengthened base from equation (9) in its Appendix A.2.2. The six-row comparison in \cref{tab:papercode} (\cref{sec:papercode}) records our interpretation, including how the unspecified index axis and latent input are implemented. We did not obtain confirmation of that interpretation from the authors. Our \dype{} port and official-code diagnostic use \href{https://github.com/guyyariv/DyPE/tree/436fe612fb7695fa2d5dc32d63390c048ba6033c}{commit \texttt{436fe612fb76}}.


\paragraph{Printed \sega{} equations}
\sega~\citep{sega2026} reads the spectrum of its input at every model call and sets per-pair gains around a reference gain $m_{\mathrm{ref}}=s^{0.08}$. In the printed equations, the energy $E_a$ of the channel-mean latent at pair $j$'s frequency bin is standardized across pairs, $z_{a,j}$, and
\begin{equation}
\begin{aligned}
m_{a,j}&=m_{\mathrm{ref}}\bigl(1-\varsigma\,u_{a,j}\bigr),
&\varsigma&=1-\mathrm{SF}(E_{\mathrm{iso}})^{1.5},\\
u_{a,j}&=\tanh z_{a,j}-\operatorname{mean}_i\tanh z_{a,i},
\end{aligned}
\label{eq:sega}
\end{equation}
where $\mathrm{SF}$ is the flatness of the radial spectrum. Low-energy pairs, usually the fast ones, thus gain attention weight in proportion to how far the spectrum is from flat. Our reading of the printed formulation differs from the release at each of the six stages in \cref{tab:papercode}. We call its parity-tested port \segaofficial{} and our implementation of the printed equations \segapaper; static variants use the mean allocation over pilot prompts and steps.

\paragraph{\dype}
For spatial axis $a$ with scale $s$ and $\kappa=\sigma^2$, let $b(r)=D\ln\bigl(64/(2\pi r)\bigr)/(2\ln\theta)$ be the pair index that completes $r$ rotations over 64 tokens and $\mu_{r_1,r_2}(j)=1-\operatorname{clip}\bigl((j-\ell)/(h-\ell),0,1\bigr)$ with $\ell=\max(\lfloor b(r_1)\rfloor,0)$ and $h=\min(\lceil b(r_2)\rceil,D/2)$. Then
\[
\begin{aligned}
\omega^{\mathrm{mix}}_j&=\tfrac{\omega_j}{s}\bigl(1-\mu_{1.25\kappa,0.75\kappa}(j)\bigr)+\omega^{\mathrm{ntk}}_j\,\mu_{1.25\kappa,0.75\kappa}(j),\\
\omega'_j&=\omega^{\mathrm{mix}}_j\bigl(1-\mu_{16\kappa,2\kappa}(j)\bigr)+\omega_j\,\mu_{16\kappa,2\kappa}(j),
\end{aligned}
\]
where $\omega^{\mathrm{ntk}}$ are the NTK-aware frequencies; the spatial cos/sin amplitude is $1+0.1\ln s$, as in the official code at the inspected revision.

\paragraph{\segaofficial}
The sensor takes the channel mean $Z$ of the projected input features on the token grid. The axis profile $E_a$ is the one-dimensional FFT power along axis $a$ averaged over the other axis and rebinned into $K_a=\max(\lfloor L_a/2\rfloor,8)$ bins. Pair $j$ reads $\ell_j=\ln E_a\bigl((1-\rho_j)(K_a-1)\bigr)$ by linear interpolation, with $\rho_j=j/(D/2-1)$ the normalized rank of its log period; $z_j$ standardizes $\ell_j$ with the sample standard deviation; and $m_{a,j}=\max\bigl(1,m_{\mathrm{ref}}[1-0.15\,\varsigma\,c_j]\bigr)$ with $c_j=\tanh(1.5z_j)-\operatorname{mean}_i\tanh(1.5z_i)$, $m_{\mathrm{ref}}=\max(s,1)^{0.08}$ and $\varsigma=1-\mathrm{SF}^{1.5}$ from the radially binned two-dimensional power. The index axis keeps gain 1.

\begin{figure}[!htbp]
\centering
\figorplaceholder[\linewidth]{figs/fA_method_zoom.pdf}{3.5in}
\caption{\textbf{Both gain controllers at component level.} (a)~The released \sega{} sensor of \cref{fig:method}b reads the channel mean of the projected input features, maps the axis profile to rotary pairs by rank and scales its tilt by the radial flatness. (b)~\constancytm{} gates a bounded log-power correction toward the time-matched native target and adds it to the \sega{} temperature. White boxes are inputs and outputs; the equations repeat \cref{alg:sega,alg:ctm} and \cref{app:impl}.}
\label{fig:method_zoom}
\end{figure}

\subsection{Per-band control laws}\label{app:laws}

Each law sets $\ln g_{a,j}=\ln\tau+\Delta_{a,j}$, a temperature plus a bounded per-band correction. The diffusion noise level supplies the bandwise noise floor.

\paragraph{\constancy}
For half-octave bands $b$, a gate compares the current input's band energy with the known noise floor, $\rho_{a,b}=\max\bigl(E^{x}_{a}[b]/(\sigma^2C)-1,0\bigr)$ and $w_{a,b}=\rho_{a,b}/(1+\rho_{a,b})$, and an error compares the previous clean estimate's band energy $\hat E_a[b]$ with a native-resolution target $P^{\mathrm{nat}}$, the mean log spectrum of the model's own $1024^2$ samples on 16 calibration prompts outside the evaluation set:
\begin{equation}
e_{a,b}=\tfrac12\bigl(P^{\mathrm{nat}}_{a}(f_b)+c_a-\ln \hat E_{a}[b]\bigr),\qquad \delta_{a,b}=\eta\,w_{a,b}\operatorname{clip}(e_{a,b},-1,1),
\label{eq:constancy}
\end{equation}
where $c_a$ anchors the target level on visible bands of 16 to 64 token wavelengths. The pair correction $\Delta_{a,j}$ interpolates $\delta$ in $\log f$, with $\tau=\tau_N$ and $\eta=1$. A natural spectrum gives $e=0$ and reduces the law to the global temperature. The target is the native spectrum of the final, nearly clean steps.

\paragraph{\constancytm}
The same law with a time-matched target: $P^{\mathrm{nat}}(f,\sigma)$ is the clean estimate's native spectrum at the same noise level, tabulated over all \N{c@steps} native steps, so the anchor no longer needs the gate. Built on the base and temperature of \segaofficial, it isolates the allocation, with \tempsb{} as its control.

\paragraph{\aliastemp{} and \segactm}
\aliastemp{} applies the log-$N$ temperature pair by pair, taking as a pair's key pool the number of keys it can resolve on the grid (\cref{app:impl}); \aliastemps{} shifts all its log gains by one constant to match \sega's reference gain in geometric mean. \segactm{} multiplies the allocation of \segaofficial{} by the time-matched correction, $\ln g=\ln g^{\textsc{sega}}+\Delta$, with $\eta\in\{0.5,1\}$ on the dev set and $\eta=1$ held out.

\paragraph{Band energies and calibration}
Band energies are densities: with an orthonormal FFT, $P(k)=\sum_c|\mathcal F(x_c-\bar x_c)(k)|^2$, $E_a(k)$ averages $P$ over the other axis, and $E_a[b]$ averages over half-octave bands $b$ (13 bands at 256 tokens per side). Unit white noise gives $E_a[b]=C=64$ in expectation, so the noise floor at level $\sigma$ is $\sigma^2C$. The gate uses the current input; the error uses the clean estimate of the previous call, $\hat x_0=x_{n-1}-\sigma_{n-1}v_\theta(x_{n-1})$. Corrections are interpolated in $\log f$ between band centers and applied to pairs with frequency above one cycle per grid, which at $4096^2$ are pairs 0 to 8 of each axis. Step 0 has no clean estimate and uses the temperature alone, and $|\ln g-\ln\tau|\le\eta$ always. The native targets come from $1024^2$ samples of the same model on 16 Pick-a-Pic prompts~\citep{kirstain2023pickapic} disjoint from the evaluation set: the late target averages the three steps with $\sigma<0.3$, and the time-matched target keeps all 28 steps and is interpolated in $\sigma$.

\paragraph{\aliastemp}
\aliastemp{} applies the log-$N$ temperature of \cref{sec:background} pair by pair. A pair of wavelength $\lambda_{a,j}$ tokens under the released base, on a grid of $L$ tokens per side ($L=64$ at native size), gets the logit gain
\[
g_{a,j}=\frac{\ln\bigl(n(\lambda_{a,j},L)\,L\bigr)}{\ln\bigl(n(\lambda_{a,j},64)\cdot 64\bigr)},\qquad n(\lambda,L)=\min\bigl(L,\max(1,L/\lambda)\bigr),
\]
where $n(\lambda,L)$, the pair's alias count, is the number of its wavelengths across the grid, bounded by 1 and $L$; the index axis takes $\tau_N$. The gains are static and content-free, with no correction and no further temperature. \aliastemps{} shifts every $\ln g$ by one constant so that the mean over all 128 dimensions equals \sega's reference log gain, $0.16\ln s$.


\section{Mechanism diagnostics}\label{app:mech}

This appendix expands the diagnostic paragraph of \cref{sec:laws} with the figure, numbers and protocols behind it.

\begin{figure}[!htbp]
\centering
\figorplaceholder[\linewidth]{figs/f7_mechanism.pdf}{2.3in}
\caption{\textbf{Dev-prompt mechanism diagnostics.} The late-profile correction starts late, and the response to a band-gain change differs across prompts. (a)~Mean $|\ln g-\ln\tau|$ over spatial rotary pairs per step: median and interquartile band over eight dev prompts, seed 0, with both targets on the \ntk{} base. In the gray region the late-profile correction is exactly zero in all eight prompts (steps 1 to 19, $\sigma\geq 0.64$). \emph{Right:} median corrections for pairs 0--11; pairs 12--27 are identically zero. (b)~Width-axis response matrices for dev prompts 0 to 2, seed 0. Each cell is the same-seed finite-difference change in final-latent log power divided by $\ln(\NmechProbeFactor)$ when one input octave receives a $\NmechProbeFactor\times$ logit gain on both spatial axes throughout sampling. The matrices share one zero-centerd color scale and are not averaged across prompts; no seed-to-seed uncertainty is estimated.}
\label{fig:mechanism}
\end{figure}

The following gate and transfer-probe measurements use development prompts; the time-matched and released-allocation summaries use logged held-out trajectories. These diagnostics do not identify causal mediators of image quality.

\paragraph{The visibility gate opens too late}
\constancy's correction is exactly zero until a band with a wavelength of 16 to 64 tokens holds more than twice the noise power ($w>\tfrac12$). On the eight dev prompts it first becomes nonzero at $\sigma=\NmechConstDevOnsetSigmaHi$ and is nonzero on all of them only by $\sigma=\NmechConstDevOnsetSigmaLo$ (\cref{fig:mechanism}a). For most of sampling, including the steps that fix the layout, \constancy{} is thus the \ntk{} base with a log-$N$ temperature: the gate is principled, but at $4096^2$ the bands it guards stay below the noise floor until late.

\paragraph{A native-resolution target improves spectral match while lowering MUSIQ}
The time-matched target opens the gate early. On the held-out \constancytm{} runs, corrections exceed $\NmechMaterialThreshold$ in log gain after the first $\NmechTmMaterialStepMed$ steps (median), and $\NmechTmNegPct\%$ of all nonzero corrections are negative: they lower the mid pairs 6 to 8 ($\NmechTmCoarseFinal$ at the last step) and, late, the finest pairs 0 to 2 ($\NmechTmFineFinal$). These corrections reduce the gains of the corresponding rotary bands. Against \tempsb, output spectral errors shrink slightly while both MUSIQ scores fall. Matching a native spectrum and maximizing a learned quality score are different objectives; the measurements do not show which individual gain changes cause the score losses.

\paragraph{\sega's allocation is mostly a fixed tilt}
Across the held-out prompts the released logit gain stays between $\NmechSegaGainMin$ and $\NmechSegaGainMax$. A fixed tilt over pairs explains $\NmechSegaStaticPct\%$ of its variance and a tilt plus a per-step profile $\NmechSegaStaticTimePct\%$. The prompt-dependent remainder, $\NmechSegaContentPct\%$ overall and $\NmechSegaLateContentPct\%$ from step $\NmechLateFirstStep$ on, is an upper bound on content dependence, because it also contains each prompt's noise draw. At the last step the finest pair has a logit gain of $\NmechSegaFinalFineGain$ and the $\NmechSegaFinalClampPairs$ slowest pairs per axis sit at the clamp of 1. A law that adapts more to content therefore has little to copy: \segactm{} raises the prompt-dependent share to $\NmechCtmContentPct\%$ and gains nothing held out. A CPU study of the printed equations finds the same structure on another latent space (\cref{app:mech}).

\paragraph{The response to gain changes is weak and prompt-dependent}
Every law assumes that raising a band's gain raises its output power. A probe on three dev prompts that scales one octave's gain by $\NmechProbeFactor$ (\cref{fig:mechanism}b) finds a clear same-octave width-axis response for one prompt and a near-zero one for the other two, off-diagonal responses spread across octaves, and a weak height-axis response. On two dev prompts, a local linear model predicts half-step band-power changes little better than no change (\cref{app:mech}). These probes question a simple band-local response model, but their small sample does not establish a universal limit on per-band control.

\paragraph{CPU study of the printed \sega{} law}
Before running any model we measured the printed allocation (\cref{eq:sega}) on forward-noised latents of the FLUX.2 autoencoder for \NmechCpuN{} DIV2K crops. A fixed tilt over pairs explained $\NmechCpuStaticPct\%$ of the variance of the modulation, tilt plus timestep $\NmechCpuStaticTimePct\%$, leaving at most $\NmechCpuContentPct\%$ for content; the flatness gate $\varsigma$ was explained to $R^2=\NmechCpuSigmaTimeRsq$ by the timestep alone and to $\NmechCpuSigmaGenericRsq$ by a generic $1/f$ field. On our held-out \krea{} runs the released gate $\varsigma$ is likewise explained to $R^2=\NmechSegaSpreadTimeRsq$ by the timestep. The channel mean that the printed sensor reads keeps $\NmechCpuChanMeanPct\%$ of the latent power. In a blur sweep at $\sigma=\NmechBlurLowSigma$ over \NmechBlurN{} test crops, the printed law's contrast between fine and coarse pairs changed by $\NmechBlurSegaChangePct\%$ between no blur and a blur of \NmechBlurMax{} tokens, while an oracle version of the constancy law changed \NmechBlurConstFold-fold. These measurements motivated the constancy design; they use the printed law, another autoencoder and forward-noised rather than sampled latents, so we use them only as context for \cref{sec:laws}.

\paragraph{Probe protocol}
The transfer probe of \cref{fig:mechanism} reruns a dev prompt with the gain of one input octave multiplied by \NmechProbeFactor{} on both spatial axes, with the same seed, and reports the finite-difference elasticity $\partial\log P_{\mathrm{out}}/\partial\log g_{\mathrm{in}}$ for six input and seven output octaves of the final latent and image. It used the \ntk{} base with the log-$N$ temperature, three dev prompts and one seed. The local-response experiment applied half-step perturbations at steps 6, 14 and 22 of two dev prompts and compared the observed change in band power with the prediction of a locally fitted linear response: the median error relative to predicting no change was $\NmechCrSeventeenErrMed$ (range $\NmechCrSeventeenErrMin$ to $\NmechCrSeventeenErrMax$; \NmechCrSeventeenBelowOne{} of \NmechCrSeventeenAxisRows{} cases below 1). It tested a single octave, below the resolution of the token grid, so it does not cover all rotary bands. On the width axis of \krea, the mean diagonal elasticity of the transfer probe is $\NmechProbeLatWDiagA$ for one prompt but $\NmechProbeLatWDiagB$ and $\NmechProbeLatWDiagC$ for the other two, against a mean off-diagonal response of $\NmechProbeLatWOffAvg$.

\paragraph{Confounded cross-checkpoint probe}
Raising the fastest octave's gain gave a width-axis elasticity of $\NmechProbeLatWTopA$ on \krea{} and $\NmechFdevLatWTop$ on \fluxdev{} for one dev prompt. Guidance and code revision also differed, so we draw no conclusion about checkpoint-dependent response from the sign change.

\paragraph{Allocation controls on the dev set}
A third dev campaign replaced the allocation of \segaofficial{} by a uniform gain with the same root-mean-square value, or kept its shape and removed the rise in mean gain caused by the clamp. On the eight dev prompts the uniform control changed MUSIQ by $\NmechCrFourteenUniformMusiq$ and the shape-only control by $+\NmechCrFourteenNormMusiq$, while a rerun of \segaofficial{} itself moved MUSIQ by $+\NmechCrFourteenRepeatMusiq$. These differences are within run-to-run variation and were not taken to held-out prompts.

\paragraph{Diagnostics not reported}
The project also measured run-to-run repeatability of the gain histories, the sensitivity of \sega's sensor to channel-mean shifts at artificial states, attention-operator approximations and native-resolution semantic consistency. None bears on the claims of this paper, and we do not report them.

\section{Compute and timing}\label{app:compute}

This appendix expands the compute paragraph of \cref{sec:protocol}: \cref{tab:cost} summarizes the kernel and GPU effects, and \cref{tab:compute} gives the cost of every retained image.

\begin{table}[!htbp]
\centering
\caption{\textbf{Kernel and GPU set cost.} Median seconds per $4096^2$ sample; the FlashAttention-2 timings come from the dev pilot. \Cref{tab:compute} expands this table. Key result: the kernel changes cost far more than the method does.}
\label{tab:cost}
\scriptsize
\setlength{\tabcolsep}{3pt}
\begin{tabular}{@{}l r r@{}}
\toprule
Setting & s/image$\downarrow$ & relative \\
\midrule
\rowcolor{papertint}GH200, cuDNN SDPA & \NcompGHcudnnSec & 1.00 \\
GH200, FlashAttention-2 & \NcompGHflashSec & \NcompFlashRatio \\
GB200, same code & \NcompBWSec & $1/\NcompBWSpeedup$ \\
\bottomrule
\end{tabular}
\end{table}

On the same GH200 node and method, forcing FlashAttention-2 makes a $4096^2$ sample take \NcompFlashRatio{} times as long as cuDNN-first SDPA (medians \NcompGHflashSec\,s and \NcompGHcudnnSec\,s), while method families on one node differ by at most \NcompFamilySpreadPct\%; one GB200 GPU is \NcompBWSpeedup{} times faster than a GH200 (\cref{tab:compute}). The forced-FlashAttention-2 comparison uses development-pilot runs and is descriptive. Timings of different methods are comparable only with the kernel and the GPU fixed.

\Cref{tab:compute} gives the sampling cost of the \NcompNimgs{} retained $4096^2$ images, \NcompGPUhours{} GPU-hours in total (\NcompGPUhoursGH{} on GH200, \NcompGPUhoursBW{} on GB200), with a median peak allocation of \NcompPeakGB\,GB. The second campaign's dev runs and the held-out \segactm{} run used one GB200 node; every other run used GH200 nodes, and every forced-FlashAttention-2 timing at $4096^2$ comes from the eight-prompt pilot. The GPU of each image was recovered from task and fleet logs, and its attention policy from the code revision pinned by its implementation hash, because the records carry a constant attention label. % Generated by paper/tables/compute_table.py. Do not edit by hand.
The SDPA column of \cref{tab:compute} gives the requested backend policy (cuDNN first: cuDNN, then FlashAttention, then memory-efficient); the dispatched kernel was not logged. The \sega{} family comprises released and paper-equation \sega{} with their static ablations; the \segactm{} family, \segactm{} and its dev-only allocation controls; the \constancy{} family, \constancy{} and its dev sweeps, \constancytm{}, the temperature-only ablations, \aliastemp{} and \aliastemps{}. Node c619-102 ran 14\% slower than the other GH200 nodes under both kernels (median power 665\,W versus 836\,W while generating), so it is listed separately; the forced-FA2 \yarn{} and \dype{} pilots ran only there. The native rows run \krea{} without RoPE changes; the first image of each worker process carries a one-off warm-up (4.1\,s with cuDNN first, 0.9\,s with FA2), which widens the cuDNN-first IQR at $1024^2$. The official \dype{} script uses the default SDPA dispatch of diffusers.
 Timings exclude model loading (\NcompModelLoadSec\,s per worker process) and PNG writing, except the single official \dype{} run.

% Generated by paper/tables/compute_table.py. Do not edit by hand.
\begin{table}[t]
\centering
\small
\caption{\textbf{Sampling cost is set by the attention kernel and the GPU, not by the method.} On the same GH200 node and method, forced FlashAttention-2 (FA2) takes 1.54 times as long as cuDNN-first SDPA (290\,s versus 188\,s per $4096^2$ image), and one GB200 GPU is 2.01 times faster than a GH200 under cuDNN-first SDPA (93\,s). Method families on one node differ by at most 0.6\%. Node c619-102 ran 14\% slower than the other GH200 nodes under both kernels (median power 665\,W versus 836\,W while generating), so it is listed separately; the forced-FA2 \yarn{} and \dype{} pilots ran only there. $n$ counts images. Seconds cover prompt encoding, 28 denoising steps and VAE decoding, without model loading or PNG writing; IQR is the interquartile range; peak is \texttt{torch.cuda.max\_memory\_allocated}; GPU-h sums the per-image seconds. The SDPA column gives the requested backend policy (cuDNN first: cuDNN, then FlashAttention, then memory-efficient); the dispatched kernel was not logged. \sega{} family: released and paper-equation \sega{} with their static ablations. \segactm{} family: \segactm{} and its dev-only allocation controls. \constancy{} family: \constancy{} and its dev sweeps, \constancytm{}, the temperature-only ablations, \aliastemp{} and \aliastemps{}. The native rows run \krea{} without RoPE changes; the first image of each worker process carries a one-off warm-up (4.1\,s with cuDNN first, 0.9\,s with FA2), which widens the cuDNN-first IQR at $1024^2$. The official \dype{} script uses the default SDPA dispatch of diffusers, and its single timing also includes writing the PNG.}
\label{tab:compute}
\begin{tabular}{@{}llrrrrrr@{}}
\toprule
SDPA & Methods & $n$ & Nodes & Median (s) & IQR (s) & Peak (GB) & GPU-h \\
\midrule
\multicolumn{8}{@{}l}{\emph{$4096^2$, GH200 (10 nodes)}} \\
FA2, forced & \ntk & 8 & 1 & 289.8 & 289.8--289.9 & 40.8 & 0.6 \\
FA2, forced & \sega{} family & 23 & 3 & 289.8 & 289.7--290.1 & 40.8 & 1.9 \\
FA2, forced & \constancy{} family & 16 & 2 & 290.4 & 289.9--290.5 & 40.8 & 1.3 \\
cuDNN first & \ntk & 160 & 6 & 187.2 & 186.2--187.5 & 40.8 & 8.3 \\
cuDNN first & \dype & 167 & 6 & 187.9 & 187.6--188.1 & 40.8 & 8.7 \\
cuDNN first & \sega{} family & 197 & 9 & 187.8 & 187.5--188.2 & 40.8 & 10.3 \\
cuDNN first & \segactm{} family & 24 & 2 & 189.4 & 184.9--189.6 & 40.8 & 1.3 \\
cuDNN first & \constancy{} family & 532 & 8 & 187.9 & 187.6--188.6 & 40.8 & 27.8 \\
\addlinespace[3pt]
\multicolumn{8}{@{}l}{\emph{$4096^2$, GH200 node c619-102 (listed separately)}} \\
FA2, forced & \yarn, \dype & 16 & 1 & 329.6 & 329.5--329.7 & 40.8 & 1.5 \\
cuDNN first & all methods & 195 & 1 & 214.6 & 214.0--215.3 & 40.8 & 11.6 \\
\addlinespace[3pt]
\multicolumn{8}{@{}l}{\emph{$4096^2$, GB200 (1 node, 4 GPUs)}} \\
cuDNN first & \segactm{} family & 204 & 1 & 93.3 & 92.5--94.3 & 40.8 & 5.3 \\
cuDNN first & \constancy{} family & 16 & 1 & 93.2 & 92.1--93.3 & 40.8 & 0.4 \\
\addlinespace[3pt]
\multicolumn{8}{@{}l}{\emph{native $1024^2$, GH200, \krea{}}} \\
FA2, forced & no RoPE change & 16 & 1 & 5.47 & 5.47--5.48 &  & \\
cuDNN first & no RoPE change & 8 & 2 & 4.94 & 4.93--5.93 & 36.3 & \\
\addlinespace[3pt]
\multicolumn{8}{@{}l}{\emph{$4096^2$, GH200, official \dype{} script}} \\
default & \dype{} (official code) & 1 & 1 & 193.3 & & & \\
\bottomrule
\end{tabular}
\end{table}


